% Lokale Grundlagen für die Newton-Beweistabellen; keine Leitfadenbeweise.
\section*{Endliche Rechnung und Differenzen}

\Nlemma{S}{Reelle endliche Summen}{%
 \begin{gathered}
 \sigma_m(\alpha u+\beta v)=\alpha\sigma_m(u)+\beta\sigma_m(v),
 \qquad \sigma_m(0)=0,\\
 \sigma_{m+1}(u)=u_0+\sigma_m(u^+),\qquad u^+_i=u_{i+1},\\
 \bigl(\forall j<r\;u_{m+j}=0\bigr)
 \ \Longrightarrow\ \sigma_{m+r}(u)=\sigma_m(u).
 \end{gathered}
}{NewtonFiniteSums}
Für die folgende Tabelle seien $L_m$ die erste Gleichung, $T_m$ die
Verschiebungsgleichung und
$Z_r=(\forall j<r\;u_{m+j}=0)\to\sigma_{m+r}(u)=\sigma_m(u)$.
Die Parameter der jeweiligen Induktion sind beliebig festgehalten.
\begin{tabproofwide}
 \proofstepwidestar[]{L_0}{\Nwhy{Summendefinition; \(0=\alpha0+\beta0\)}}
 \proofstepwidestar[2]{L_m}{\rA}
 \proofstepwidestar[2]{\begin{aligned}[t]
 \sigma_{m+1}(\alpha u+\beta v)
 &=\sigma_m(\alpha u+\beta v)+\alpha u_m+\beta v_m\\
 &=\alpha(\sigma_m(u)+u_m)+\beta(\sigma_m(v)+v_m)
 \end{aligned}}{\Nwhy{Summendefinition; 2; Körperrechnung}}
 \proofstepwidestar[2]{L_{m+1}}{\Nwhy{3; Summendefinition}}
 \proofstepwidestar[]{\forall m\in\mathbb N\;L_m}{\NI{1;2--4}}
 \proofstepwidestar[]{\sigma_m(0)=0}{\Nwhy{5 mit \(\alpha=\beta=0\); Körperrechnung}}
 \proofstepwidestar[]{T_0}{\Nwhy{\(\sigma_1(u)=u_0=u_0+\sigma_0(u^+)\)}}
 \proofstepwidestar[8]{T_m}{\rA}
 \proofstepwidestar[8]{\begin{aligned}[t]
 \sigma_{m+2}(u)&=\sigma_{m+1}(u)+u_{m+1}\\
 &=u_0+\sigma_m(u^+)+u^+_m
 =u_0+\sigma_{m+1}(u^+)
 \end{aligned}}{\Nwhy{Summendefinition; 8; Def. \(u^+\)}}
 \proofstepwidestar[]{\forall m\in\mathbb N\;T_m}{\NI{7;8--9}}
 \proofstepwidestar[]{Z_0}{\Nwhy{\(\sigma_{m+0}(u)=\sigma_m(u)\); \(\to I\)}}
 \proofstepwidestar[12]{Z_r}{\rA}
 \proofstepwidestar[13]{\forall j<r+1\;u_{m+j}=0}{\rA}
 \proofstepwidestar[13]{(\forall j<r\;u_{m+j}=0)\land u_{m+r}=0}{\Nwhy{13; \(j<r\Rightarrow j<r+1\), \(r<r+1\)}}
 \proofstepwidestar[12,13]{\sigma_{m+r}(u)=\sigma_m(u)}{\rRE{12,\rAEa{14}}}
 \proofstepwidestar[12,13]{\sigma_{m+r+1}(u)=\sigma_{m+r}(u)+u_{m+r}=\sigma_m(u)}{\Nwhy{Summendefinition; 14,15}}
 \proofstepwidestar[12]{Z_{r+1}}{\rRI{13,16}}
 \proofstepwidestar[]{\forall r\in\mathbb N\;Z_r}{\NI{11;12--17}}
\end{tabproofwide}
Die Nullverlängerung gilt insbesondere für die Nullfolge. Mit $m=0$
besagt sie auch: Sind alle tatsächlich summierten Glieder null, so ist
die Summe null. Die Verschiebung ist die benötigte Indexumbenennung;
allgemeine Umordnungen unendlicher Summen werden nicht verwendet.

\Needspace{14\baselineskip}
\Nlemma{C}{Koeffizientenrechnung}{%
 \begin{gathered}
 \alpha P_c+\beta P_e=P_{\alpha c+\beta e},\qquad
 (\alpha X+\beta)P_c=P_b,\quad b_i=\alpha c_{i-1}+\beta c_i,\\
 d\le D\ \Longrightarrow\ \mathcal P_d\subseteq\mathcal P_D,\\
 p,q\in\mathcal P_d\ \Longrightarrow\
 \alpha p+\beta q\in\mathcal P_d\ \land\
 (\alpha X+\beta)p\in\mathcal P_{d+1},\\
 \forall i<m\;(p_i\in\mathcal P_d)
 \ \Longrightarrow\ \sum_{i<m}p_i\in\mathcal P_d.
 \end{gathered}
}{NewtonCoefficientCalculus}
In der ersten Zeile sind $c,e\in\mathbb R^{d+1}$,
$b\in\mathbb R^{d+2}$ und $c_i=0$ außerhalb von $0\le i\le d$.
Die Tupelargumente der Tabelle gelten für $d\ge0$; die eigenen Nullfolgenzeilen behandeln $d=-1$.
Bei $D\ge d$ bezeichnet $\widetilde c$ die bis $D$ mit Nullen
verlängerte Koeffizientenliste. Die letzte Aussage wird mit
$F_m=(\forall i<m\;p_i\in\mathcal P_d)\to\sum_{i<m}p_i\in\mathcal P_d$
abgekürzt.
\begin{tabproofwide}
 \proofstepwidestar[]{\alpha P_c+\beta P_e
 =\sum_{i=0}^d(\alpha c_i+\beta e_i)u_i=P_{\alpha c+\beta e}}
 {\Nwhy{Def. \(P_c\); \Nref{S}; punktweise Auswertung}}
 \proofstepwidestar[]{\begin{aligned}[t]
 (\alpha X+\beta)P_c
 &=\alpha\sum_{i=0}^d c_i u_{i+1}
   +\beta\sum_{i=0}^d c_i u_i\\
 &=\sum_{i=0}^{d+1}(\alpha c_{i-1}+\beta c_i)u_i=P_b
 \end{aligned}}{\Nwhy{\Nref{S}; \(Xu_i=u_{i+1}\); Nullverlängerung}}
 \proofstepwidestar[]{0\in\mathcal P_d}{\Nwhy{Nullkoeffizienten; \Nref{S}; für \(d=-1\) Definition}}
 \proofstepwidestar[4]{p,q\in\mathcal P_d}{\rA}
 \proofstepwidestar[4]{\exists c,e\in\mathbb R^{d+1}\;(p=P_c\land q=P_e)}{\Nwhy{Def. \(\mathcal P_d\), zunächst \(d\ge0\)}}
 \proofstepwidestar[6]{c,e\in\mathbb R^{d+1}\land p=P_c\land q=P_e}{\rA}
 \proofstepwidestar[6]{\alpha p+\beta q\in\mathcal P_d}{\Nwhy{1; 6; Def. \(\mathcal P_d\)}}
 \proofstepwidestar[6]{(\alpha X+\beta)p\in\mathcal P_{d+1}}{\Nwhy{2; 6; Def. \(\mathcal P_{d+1}\)}}
 \proofstepwidestar[6]{\alpha p+\beta q\in\mathcal P_d\land
 (\alpha X+\beta)p\in\mathcal P_{d+1}}{\rAI{7,8}}
 \proofstepwidestar[4]{\alpha p+\beta q\in\mathcal P_d\land
 (\alpha X+\beta)p\in\mathcal P_{d+1}}{\rEE{5,6,9}}
 \proofstepwidestar[]{p,q\in\mathcal P_d\to
 \alpha p+\beta q\in\mathcal P_d\land
 (\alpha X+\beta)p\in\mathcal P_{d+1}}{\rRI{4,10}}
 \proofstepwidestar[]{\begin{aligned}[t]
 p,q\in\mathcal P_{-1}\to{}&\alpha p+\beta q\in\mathcal P_{-1}\\[-2pt]
 &\land\ (\alpha X+\beta)p\in\mathcal P_0
 \end{aligned}}{\Nwhy{\(p=q=0\); 3; Körperrechnung}}
 \proofstepwidestar[13]{d\le D\land p\in\mathcal P_d}{\rA}
 \proofstepwidestar[13]{\exists c\in\mathbb R^{d+1}\;(p=P_c)}{\Nwhy{13; Def. \(\mathcal P_d\), zunächst \(d\ge0\)}}
 \proofstepwidestar[15]{c\in\mathbb R^{d+1}\land p=P_c}{\rA}
 \proofstepwidestar[13,15]{p\in\mathcal P_D}
 {\Nwhy{\Nref{S}, Nullverlängerung; 13,15; Def. \(\mathcal P_D\)}}
 \proofstepwidestar[13]{p\in\mathcal P_D}{\rEE{14,15,16}}
 \proofstepwidestar[]{d\le D\land p\in\mathcal P_d\to p\in\mathcal P_D}{\rRI{13,17}}
 \proofstepwidestar[]{\mathcal P_{-1}\subseteq\mathcal P_D}{\Nwhy{\(\mathcal P_{-1}=\{0\}\); 3, auch für \(D=-1\)}}
 \proofstepwidestar[]{F_0}{\Nwhy{Leere Summe; 3; \(\to I\)}}
 \proofstepwidestar[21]{F_m}{\rA}
 \proofstepwidestar[22]{\forall i<m+1\;(p_i\in\mathcal P_d)}{\rA}
 \proofstepwidestar[22]{(\forall i<m\;p_i\in\mathcal P_d)\land p_m\in\mathcal P_d}{\Nwhy{22; Indexeinschränkung}}
 \proofstepwidestar[21,22]{\sum_{i<m}p_i\in\mathcal P_d}{\rRE{21,\rAEa{23}}}
 \proofstepwidestar[21,22]{\sum_{i<m+1}p_i\in\mathcal P_d}{\Nwhy{Summenrekursion; 11--12 mit \(\alpha=\beta=1\); 23,24}}
 \proofstepwidestar[21]{F_{m+1}}{\rRI{22,25}}
 \proofstepwidestar[]{\forall m\in\mathbb N\;F_m}{\NI{20;21--26}}
\end{tabproofwide}
Die Koeffizientenlisten werden ausschließlich als Darstellungen benutzt;
ihre Eindeutigkeit ist hier keine Voraussetzung.

\Nlemma{P}{Pascal-Identität mit Randfällen}{%
 \binom{n+1}{k}=\binom nk+\binom n{k-1}
 \qquad(n,k\in\mathbb N),\qquad \binom n{-1}=0.
}{NewtonPascal}
\begin{tabproofwide}
 \proofstepwidestar[1]{1\le k\le n}{\rA}
 \proofstepwidestar[1]{\begin{aligned}[t]
 \binom nk+\binom n{k-1}
 &=\frac{n!(n-k+1)}{k!(n-k+1)!}
   +\frac{n!k}{k!(n-k+1)!}\\
 &=\frac{n!(n+1)}{k!(n+1-k)!}
 =\binom{n+1}k
 \end{aligned}}{\Nwhy{Binomialdefinition; \NF; Körperrechnung}}
 \proofstepwidestar[]{1\le k\le n\to
 \binom{n+1}k=\binom nk+\binom n{k-1}}
 {\Nwhy{\(\to I\), 1,2; Symmetrie}}
 \proofstepwidestar[]{\binom{n+1}0=1=\binom n0+\binom n{-1}}
 {\Nwhy{Definition; \NF; \(0!=1\)}}
 \proofstepwidestar[]{\binom{n+1}{n+1}=1=\binom n{n+1}+\binom nn}
 {\Nwhy{Definition; \NF}}
 \proofstepwidestar[6]{k>n+1}{\rA}
 \proofstepwidestar[6]{\binom{n+1}k=0=\binom nk+\binom n{k-1}}
 {\Nwhy{Definition; \(k>n\), \(k-1>n\)}}
 \proofstepwidestar[]{k>n+1\to
 \binom{n+1}k=\binom nk+\binom n{k-1}}{\rRI{6,7}}
 \proofstepwidestar[]{k=0\ \lor\ 1\le k\le n\ \lor\ k=n+1\ \lor\ k>n+1}
 {\Nwhy{Natürliche Ordnung; vollständige Fälle}}
 \proofstepwidestar[]{\binom{n+1}k=\binom nk+\binom n{k-1}}
 {\Nwhy{3--5,8,9; Fallunterscheidung}}
\end{tabproofwide}
Für $n=0$ ist der innere Fall leer; die Zeilen 4, 5 und 8 liefern
unverändert sämtliche Werte von $k$.

\Nlemma{D}{Lineare Differenzen und Iteration}{%
 \begin{gathered}
 \Delta^r(\alpha a+\beta b)=\alpha\Delta^ra+\beta\Delta^rb,
 \qquad \Delta^r0=0,\\
 \Delta^{r+s}a=\Delta^r(\Delta^sa),\\
 \Delta^r\left(\sum_{i<m}\alpha_i a^{(i)}\right)
 =\sum_{i<m}\alpha_i\Delta^ra^{(i)}.
 \end{gathered}
}{NewtonDifferenceCalculus}
Die endliche Familie erfüllt $\alpha_i\in\mathbb R$ und $a^{(i)}\in\mathbb R^{\mathbb N}$ für $i<m$. Seien $L_r$ die Linearitätsgleichung, $K_r$ die Kompositionsgleichung
bei festem $s$ und $M_m$ die Gleichung für die endliche Summe.
\begin{tabproofwide}
 \proofstepwidestar[]{\begin{aligned}[t]
 (\Delta(\alpha a+\beta b))_n
 &=(\alpha a_{n+1}+\beta b_{n+1})-(\alpha a_n+\beta b_n)\\
 &=\alpha(\Delta a)_n+\beta(\Delta b)_n
 \end{aligned}}{\Nwhy{Differenzendefinition; Körperrechnung}}
 \proofstepwidestar[]{\Delta(\alpha a+\beta b)=\alpha\Delta a+\beta\Delta b}
 {\Nwhy{1; Funktionenextensionalität}}
 \proofstepwidestar[]{L_0}{\Nwhy{\(\Delta^0a=a\)}}
 \proofstepwidestar[4]{L_r}{\rA}
 \proofstepwidestar[4]{\begin{aligned}[t]
 \Delta^{r+1}(\alpha a+\beta b)
 &=\Delta(\alpha\Delta^ra+\beta\Delta^rb)\\
 &=\alpha\Delta^{r+1}a+\beta\Delta^{r+1}b
 \end{aligned}}{\Nwhy{Differenzenrekursion; 2,4}}
 \proofstepwidestar[]{\forall r\in\mathbb N\;L_r}{\NI{3;4--5}}
 \proofstepwidestar[]{\Delta^r0=0}{\Nwhy{6 mit \(\alpha=\beta=0\); Körperrechnung}}
 \proofstepwidestar[]{K_0}{\Nwhy{\(\Delta^{0+s}a=\Delta^sa=\Delta^0(\Delta^sa)\)}}
 \proofstepwidestar[9]{K_r}{\rA}
 \proofstepwidestar[9]{\Delta^{r+1+s}a
 =\Delta(\Delta^r(\Delta^sa))=\Delta^{r+1}(\Delta^sa)}
 {\Nwhy{Differenzenrekursion; 9; \((r+1)+s=(r+s)+1\)}}
 \proofstepwidestar[]{\forall r\in\mathbb N\;K_r}{\NI{8;9--10}}
 \proofstepwidestar[]{M_0}{\Nwhy{Leere Summe; 7}}
 \proofstepwidestar[13]{M_m}{\rA}
 \proofstepwidestar[13]{\begin{aligned}[t]
 \Delta^r\left(\sum_{i<m+1}\alpha_i a^{(i)}\right)
 &=\Delta^r\left(\sum_{i<m}\alpha_i a^{(i)}\right)
   +\alpha_m\Delta^ra^{(m)}\\
 &=\sum_{i<m+1}\alpha_i\Delta^ra^{(i)}
 \end{aligned}}{\Nwhy{Summenrekursion; 6,13}}
 \proofstepwidestar[]{\forall m\in\mathbb N\;M_m}{\NI{12;13--14}}
\end{tabproofwide}

\Nlemma{B}{Binomialfolgen sind polynomiale Folgen}{%
 B_k(n)=\binom nk,\qquad B_k\in\mathcal P_k,
 \qquad B_{k+1}=\frac{X-k}{k+1}B_k
 \quad(n,k\in\mathbb N).
}{NewtonBinomialPolynomials}
Setze $J_k=(\forall n\in\mathbb N\;B_k(n)=\binom nk)
\land B_k\in\mathcal P_k$.
\begin{tabproofwide}
 \proofstepwidestar[]{B_{k+1}=\frac{X-k}{k+1}B_k}
 {\Nwhy{Produktrekursion; \NF; Körperrechnung; Extensionalität}}
 \proofstepwidestar[2]{k<n}{\rA}
 \proofstepwidestar[2]{\begin{aligned}[t]
 \binom n{k+1}
 &=\frac{n!}{(k+1)!(n-k-1)!}\\
 &=\frac{n-k}{k+1}\frac{n!}{k!(n-k)!}
 =\frac{n-k}{k+1}\binom nk
 \end{aligned}}{\Nwhy{Binomialdefinition; \NF; Körperrechnung}}
 \proofstepwidestar[]{k<n\to\binom n{k+1}=\frac{n-k}{k+1}\binom nk}{\rRI{2,3}}
 \proofstepwidestar[]{k=n\to\binom n{k+1}=0=\frac{n-k}{k+1}\binom nk}
 {\Nwhy{Binomialdefinition; \(n-k=0\)}}
 \proofstepwidestar[]{k>n\to\binom n{k+1}=0=\frac{n-k}{k+1}\binom nk}
 {\Nwhy{Binomialdefinition; \(\binom nk=0\)}}
 \proofstepwidestar[]{\binom n{k+1}=\frac{n-k}{k+1}\binom nk}
 {\Nwhy{4--6; \(k<n\lor k=n\lor k>n\)}}
 \proofstepwidestar[]{J_0}
 {\Nwhy{\(B_0=1\in\mathcal P_0\), \(\binom n0=1\); \NF}}
 \proofstepwidestar[9]{J_k}{\rA}
 \proofstepwidestar[9]{\forall n\in\mathbb N\;
 B_{k+1}(n)=\frac{n-k}{k+1}B_k(n)=\binom n{k+1}}
 {\Nwhy{1,7,9; \(\forall E,\forall I\)}}
 \proofstepwidestar[9]{B_{k+1}\in\mathcal P_{k+1}}
 {\Nwhy{1,9; \Nref{C} mit \(\alpha=1/(k+1)\), \(\beta=-k/(k+1)\)}}
 \proofstepwidestar[9]{J_{k+1}}{\Nwhy{10,11; Transitivität; \(\land I\)}}
 \proofstepwidestar[]{\forall k\in\mathbb N\;J_k}{\NI{8;9--12}}
\end{tabproofwide}
